\chapter{Ergebnisse und Diskussion}
\label{chap:ergebnisse}
%
Dieses Kapitel präsentiert die experimentellen Ergebnisse und interpretiert sie im Kontext der Forschungsfrage.
Die Darstellung folgt der Phasenstruktur des Experiments: Für jede Phase werden zunächst die Ergebnisse berichtet, anschließend im Hinblick auf die in der Methodik formulierten Erwartungen diskutiert und schließlich Implikationen für die nachfolgenden Phasen abgeleitet.
%
\section{Phase~1: Charakterisierung des Float-Raums}
\label{sec:ergebnisse_phase1}
%
Phase~1 analysiert den unkomprimierten Float-Raum, um Erwartungen an das Quantisierungsverhalten zu formulieren.
Die Ergebnisse zeigen, dass der von \texttt{bge-large-en-v1.5} erzeugte Embedding-Raum Eigenschaften besitzt, die Quantisierung systematisch begünstigen.
%
\subsection{Normverteilung}
%
Tabelle~\ref{tab:phase1_norms} fasst die Normstatistiken über alle drei Datensätze zusammen.
%
\begin{table}[htbp]
  \centering
  \caption{L2-Normverteilung der Float32-Embeddings. CV bezeichnet den Variationskoeffizienten $\sigma / \mu$.}
  \label{tab:phase1_norms}
  \begin{tabular}{lcccc}
    \toprule
    Datensatz & $\mu$ & $\sigma$ & Bereich & CV \\
    \midrule
    SciFact     & 1{,}000\,000 & $2{,}93 \times 10^{-8}$ & $[1{,}000\,000;\; 1{,}000\,000]$ & $2{,}93 \times 10^{-8}$ \\
    FiQA        & 1{,}000\,000 & $2{,}96 \times 10^{-8}$ & $[1{,}000\,000;\; 1{,}000\,000]$ & $2{,}96 \times 10^{-8}$ \\
    TREC-COVID  & 1{,}000\,000 & $3{,}03 \times 10^{-8}$ & $[1{,}000\,000;\; 1{,}000\,000]$ & $3{,}03 \times 10^{-8}$ \\
    \bottomrule
  \end{tabular}
\end{table}
%
Die Variationskoeffizienten von $\text{CV} \approx 3 \times 10^{-8}$ bestätigen, dass die Vektoren exakt auf der Einheitshypersphäre liegen.
Die beobachteten Abweichungen liegen in der Größenordnung der Float32-Maschinengenauigkeit ($\epsilon_{\text{float32}} \approx 1{,}2 \times 10^{-7}$) und sind damit numerisches Rauschen, keine inhaltliche Variation.
%
\paragraph{Interpretation.}
Dieses Ergebnis hat eine direkte Konsequenz für die Binarisierung.
Da die Vorzeichenfunktion die Norminformation vollständig verwirft (vgl.\ Abschnitt~\ref{sec:methodik_phase1}), hängt der Informationsverlust davon ab, wie viel Information in der Norm kodiert ist.
Bei perfekter Normierung ist der Informationsgehalt der Norm null.
Die Binarisierung verliert damit ausschließlich die Magnitudeninformation \emph{innerhalb} jeder Dimension, nicht die Gesamt-Norminformation.
Dies ist eine günstige Voraussetzung, die nicht für alle Embedding-Modelle gelten muss.
%
\subsection{Dimensionsstatistiken}
%
Tabelle~\ref{tab:phase1_dimstats} zeigt die aggregierten Dimensionsstatistiken.
Die zugehörige Verteilung über alle 1024 Dimensionen ist in Abbildung~\ref{fig:dimension_distribution} dargestellt.
%
\begin{table}[htbp]
  \centering
  \caption{Dimensionsstatistiken über alle 1024 Dimensionen. Die Schwellwerte $|\text{skew}| > 1{,}0$ und $\text{kurtosis} > 3{,}0$ (Fisher-Kurtosis, Normalverteilung $= 0$) markieren problematische Dimensionen für Quantisierung.}
  \label{tab:phase1_dimstats}
  \begin{tabular}{lccccc}
    \toprule
    Datensatz & $\overline{|\text{skew}|}$ & $\max |\text{skew}|$ & $\overline{\text{kurt}}$ & $\max(\text{kurt})$ & $n_{\text{probl.}}$ \\
    \midrule
    SciFact     & 0{,}074 & 0{,}377 & 0{,}003 & 0{,}662 & 0 / 1024 \\
    FiQA        & 0{,}081 & 0{,}471 & 0{,}044 & 0{,}665 & 0 / 1024 \\
    TREC-COVID  & 0{,}085 & 0{,}366 & 0{,}022 & 0{,}539 & 0 / 1024 \\
    \bottomrule
  \end{tabular}
\end{table}
%
\begin{figure}[htbp]
  \centering
  \includegraphics[width=\textwidth]{dimension_distribution.pdf}
  \caption{Verteilung der Schiefe und Kurtosis über alle 1024 Dimensionen für die drei Datensätze. Die gestrichelten Linien markieren die Schwellwerte $|\text{skew}| = 1{,}0$ und $\text{kurtosis} = 3{,}0$, ab denen Quantisierungsprobleme zu erwarten wären.}
  \label{fig:dimension_distribution}
\end{figure}
%
Die Ergebnisse sind über alle drei Datensätze bemerkenswert konsistent.
Keine einzige der 3072 untersuchten Dimensionen ($3 \times 1024$) überschreitet die Schwellwerte für problematische Schiefe oder Kurtosis.
Die maximale Schiefe liegt bei 0{,}471 (FiQA), also weniger als der Hälfte des Schwellwerts.
Die maximale Kurtosis beträgt 0{,}665 (FiQA), was einer nur leicht leptokurtischen Verteilung entspricht, die weit entfernt von problematischem Heavy-Tail-Verhalten ist.
%
\paragraph{Interpretation.}
Die nahezu symmetrischen Verteilungen um Null haben zwei konkrete Implikationen.
%
Erstens liegt der Vorzeichenschwellwert $x_i = 0$ der Binarisierung nahe am informationstheoretisch optimalen Punkt.
Bei einer symmetrischen Verteilung teilt der Schwellwert~0 die Wahrscheinlichkeitsmasse annähernd in zwei gleich große Hälften, was die Entropie des resultierenden Bits maximiert.
Die in der Methodik formulierte Erwartung, dass schiefe Dimensionen bei der Binarisierung mehr Fehler verursachen sollten, lässt sich hier nicht prüfen, da keine solchen Dimensionen existieren.
Das Fehlen schiefer Dimensionen ist jedoch selbst ein Ergebnis: Es zeigt, dass \texttt{bge-large-en-v1.5} einen für Binarisierung günstigen Raum erzeugt.
%
Zweitens bedeutet die niedrige Kurtosis, dass keine Dimension Ausreißer aufweist, die bei uniformer Quantisierung mit wenigen Stufen (etwa 4~Stufen bei 2-Bit) zu großen Fehlern führen würden.
Die uniforme Stufenverteilung von TurboQuant ist damit für diesen Raum gut geeignet.
%
Die Konsistenz über drei unterschiedliche Domänen (Naturwissenschaft, Finanz, Biomedizin) deutet darauf hin, dass diese Eigenschaften vom Modell und nicht vom Datensatz bestimmt werden.
Dies ist plausibel, da das Training über diverse Korpora hinweg eine symmetrische Nutzung des Raums erzwingt.
%
\subsection{Intrinsische Dimensionalität}
\label{sec:ergebnisse_intdim}
%
Tabelle~\ref{tab:phase1_intdim} stellt die Schätzungen der intrinsischen Dimensionalität gegenüber.
Das zugehörige PCA-Varianzspektrum und die kumulative erklärte Varianz sind in den Abbildungen~\ref{fig:pca_variance_spectrum} und~\ref{fig:pca_cumulative_variance} dargestellt.
%
\begin{table}[htbp]
  \centering
  \caption{Intrinsische Dimensionalität: TwoNN (lokale Mannigfaltigkeitsstruktur) und PCA (globale lineare Dimensionalität, 95\% erklärte Varianz). Ergänzend die kumulative Varianz bei ausgewählten Komponentenzahlen.}
  \label{tab:phase1_intdim}
  \begin{tabular}{lccccc}
    \toprule
    Datensatz & TwoNN & PCA$_{95\%}$ & $V(50)$ & $V(100)$ & $V(256)$ \\
    \midrule
    SciFact     & 13{,}2 & 390 & 50{,}0\% & 65{,}2\% & 87{,}4\% \\
    FiQA        & 22{,}1 & 395 & 52{,}7\% & 67{,}0\% & 87{,}8\% \\
    TREC-COVID  & 11{,}1 & 417 & 50{,}0\% & 63{,}8\% & 85{,}4\% \\
    \bottomrule
  \end{tabular}
\end{table}
%
\begin{figure}[htbp]
  \centering
  \includegraphics[width=\textwidth]{pca_variance_spectrum.pdf}
  \caption{PCA-Varianzspektrum: Erklärte Varianz pro Hauptkomponente (absteigende Reihenfolge). Der schnelle Abfall der ersten Komponenten gefolgt von einem langen, flachen Schwanz zeigt die typische Varianzkonzentration in hochdimensionalen Embedding-Räumen.}
  \label{fig:pca_variance_spectrum}
\end{figure}
%
\begin{figure}[htbp]
  \centering
  \includegraphics[width=\textwidth]{pca_cumulative_variance.pdf}
  \caption{Kumulative erklärte Varianz als Funktion der Komponentenzahl. Die horizontale Linie markiert 95\%. Die vertikalen Markierungen zeigen die evaluierten PCA-Dimensionen $d \in \{64, 128, 256, 384, 512, 768, 1024\}$.}
  \label{fig:pca_cumulative_variance}
\end{figure}
%
\paragraph{Diskrepanz zwischen lokaler und globaler Schätzung.}
Die auffälligste Beobachtung ist die starke Divergenz zwischen den beiden Schätzern.
TwoNN schätzt die lokale Mannigfaltigkeitsdimension auf 11--22, während PCA~390--417 Komponenten für 95\% Varianz benötigt.
Diese Diskrepanz beträgt mehr als eine Größenordnung und ist kein Artefakt, sondern inhaltlich interpretierbar.
%
TwoNN misst die lokale Freiheitsgrade der Datenmannigfaltigkeit: Wie viele unabhängige Richtungen existieren in der unmittelbaren Umgebung eines Punktes?
Der niedrige Wert (11--22) zeigt, dass die Dokumente lokal auf einer niedrigdimensionalen Mannigfaltigkeit liegen.
Semantisch ähnliche Dokumente unterscheiden sich nur entlang weniger Achsen.
%
PCA misst hingegen die globale lineare Dimensionalität: Wie viele orthogonale Richtungen sind nötig, um die Gesamtstreuung aller Punkte zu erfassen?
Der hohe Wert (390--417) zeigt, dass die Mannigfaltigkeit zwar lokal niedrigdimensional ist, aber global viele verschiedene Richtungen einnimmt.
Dies entspricht dem Bild einer gekrümmten Mannigfaltigkeit, die zwar lokal flach ist, aber in einen hochdimensionalen Raum eingebettet ist und dort viele verschiedene Orientierungen annimmt.
%
\paragraph{Implikationen für die Dimensionsreduktion.}
Aus den PCA-Ergebnissen lassen sich quantitative Vorhersagen für die nachfolgenden Phasen ableiten:
\begin{itemize}
  \item \textbf{$d = 384$:} Mit 384 Komponenten werden je nach Datensatz 93--94\% der Varianz erhalten (interpoliert aus den Werten für $V(256)$ und $V(\text{PCA}_{95\%})$).
Eine moderate Qualitätsreduktion ist zu erwarten, aber keine katastrophale.
  \item \textbf{$d = 256$:} Bei 256 Komponenten verbleiben 85--88\% der Varianz.
Hier beginnt der Bereich, in dem substantielle Information verloren geht.
  \item \textbf{$d = 64$:} Bei 64 Komponenten verbleiben nur ca. 40\% der Varianz.
Ein massiver Qualitätsverlust ist zu erwarten.
\end{itemize}
%
\paragraph{Implikationen für die Quantisierung.}
Die niedrige lokale Dimensionalität (TwoNN $\approx$ 11--22) hat eine weitere Konsequenz.
Sie bedeutet, dass die lokale Nachbarschaftsstruktur, die für Retrieval entscheidend ist, durch relativ wenige effektive Dimensionen bestimmt wird.
Quantisierung, die jede Dimension unabhängig behandelt, kann daher einen substantiellen Teil der Bitbreite auf redundante Dimensionen verwenden, ohne die relevante lokale Struktur zu beschädigen.
Dies legt nahe, dass selbst aggressive Quantisierung (2-Bit, 1-Bit) die \emph{lokale} Nachbarschaftsstruktur besser erhalten könnte als die \emph{globale} Distanzstruktur, da erstere von weniger effektiven Dimensionen abhängt.
Diese Hypothese wird in Phase~3 empirisch geprüft.
%
\subsection{Zusammenfassung und Vorhersagen}
\label{sec:ergebnisse_phase1_vorhersagen}
%
Die Charakterisierung ergibt ein konsistentes Bild: Der Embedding-Raum von \texttt{bge-large-en-v1.5} besitzt drei Eigenschaften, die Quantisierung systematisch begünstigen:
\begin{enumerate}
  \item \textbf{Perfekte Normierung} -- die Binarisierung verliert keine Norminformation.
  \item \textbf{Symmetrische, leicht-schwänzige Dimensionen} -- der Vorzeichenschwellwert liegt optimal, und uniforme Quantisierung verursacht keine Ausreißer-induzierten Fehler.
  \item \textbf{Hohe Redundanz} -- die intrinsische Dimensionalität liegt weit unter der nominalen, sodass Kompression über beide Achsen (Bittiefe und Dimensionszahl) möglich ist.
\end{enumerate}
%
Daraus ergeben sich spezifische, testbare Vorhersagen für die nachfolgenden Phasen:
\begin{itemize}
  \item Selbst die naive Binarisierung sollte aufgrund der günstigen Raumeigenschaften eine nicht-triviale Distanz- und Nachbarschaftserhaltung zeigen.
  \item Der Vorteil von TurboQuant gegenüber naiver Quantisierung sollte moderat ausfallen, da die Varianzumverteilung durch Rotation in einem bereits gleichmäßig genutzten Raum weniger Verbesserungspotenzial hat.
  \item PCA-Reduktion auf $d \geq 384$ sollte bei allen Quantisierungsstufen nur marginale Qualitätsverluste verursachen.
  \item Die stärksten Qualitätseinbrüche werden bei $d \leq 128$ erwartet, wo mehr als 30\% der Varianz verloren gehen.
\end{itemize}
%
Ob diese Vorhersagen zutreffen, wird in den Phasen~2 bis~4 empirisch überprüft.